% Auto-generated Table 1
% Source: paper_analysis_corpus.csv descriptive.json
\begin{table}[ht]
\centering
\caption{Characteristics of the Three Cyber Conflict Databases}
\label{tab:db-comparison}
\begin{tabular}{lp{3.5cm}p{3.5cm}p{3.5cm}}
\toprule
Characteristic & CISSM & CSIS & CFR \\
\midrule
Entries & 335 & 81 & 26 \\
Time range & Feb 2022--Dec 2024 & Jan 2022--Dec 2024 & Feb 2022--Dec 2022 \\
Coding unit & Cyber operation & Significant cyber incident & Policy-relevant event \\
Event taxonomy & Disruptive/Exploitive/ Mixed/Undetermined & Narrative descriptions\textsuperscript{a} & Espionage/Sabotage/ DDoS/Data destruction \\
Attribution method & OSINT self-attestation + researcher coding & State/official attribution & Media/law enforcement \\
Data access & Closed (Georgetown) & Public (CSIS RSI) & Public (CFR Cyber Tracker) \\
\bottomrule
\end{tabular}
\vspace{4pt}
\raggedright\footnotesize{\textsuperscript{a}CSIS entries use descriptive narratives instead of a structured event-type taxonomy. Excluded from event-type analyses (see \S4).}
\end{table}
